\lyseis{Β3.10}{Τετράπλευρα, παραλληλόγραμμα, τραπέζια}
\sol{1} α) Σωστό. β) Λάθος. γ) Σωστό. δ) Σωστό. ε) Λάθος (δεν είναι γενικά ίσες). στ) Σωστό.
\sol{2} Ελεύθερη κατασκευή. Πρέπει να βρεις $\mathrm{ΟΑ}=\mathrm{ΟΓ}$ (και $\mathrm{ΟΒ}=\mathrm{ΟΔ}$), επειδή οι διαγώνιοι διχοτομούνται, και το ύψος να είναι κάθετο στη βάση.
\sol{3} $\mathrm{ΓΔ}=\mathrm{ΑΒ}=9$ cm, $\mathrm{ΑΔ}=\mathrm{ΒΓ}=5$ cm. Περίμετρος $2\cdot(9+5)=28$ cm.
\sol{4} $\hat{\mathrm{Γ}}=72^\circ$ και $\hat{\mathrm{Β}}=\hat{\mathrm{Δ}}=180^\circ-72^\circ=108^\circ$.
\sol{5} Οι διαγώνιοι διχοτομούνται: $\mathrm{ΟΑ}=\mathrm{ΟΓ}=5$ cm και $\mathrm{ΟΒ}=\mathrm{ΟΔ}=4$ cm.
\sol{6} $\hat{\mathrm{Γ}}=68^\circ$ (γωνίες της ίδιας βάσης), $\hat{\mathrm{Α}}=\hat{\mathrm{Β}}=180^\circ-68^\circ=112^\circ$.
\sol{7} Με τέμνουσα την $\mathrm{ΑΔ}$: $\hat{\mathrm{Δ}}=180^\circ-95^\circ=85^\circ$. Με τέμνουσα την $\mathrm{ΒΓ}$: $\hat{\mathrm{Γ}}=180^\circ-110^\circ=70^\circ$.
\sol{8} Η περίμετρος είναι $2\cdot(30+20)=100$ cm, δηλαδή $1$ μ. ταινία.
\sol{9} Οι γωνίες στη μεγάλη βάση είναι $70^\circ$ και $70^\circ$, και στη μικρή βάση $180^\circ-70^\circ=110^\circ$ και $110^\circ$. Η περίμετρος είναι $40+60+25+25=150$ cm.
\sol{10} Οι γωνίες είναι $60^\circ,120^\circ,60^\circ,120^\circ$. Τα μήκη των πλευρών δεν αλλάζουν, μόνο οι γωνίες.
\sol{11} Αν η μικρή πλευρά είναι $x$, η άλλη είναι $x+3$ και $2\cdot(x+x+3)=46$, άρα $2x+3=23$ και $x=10$. Οι πλευρές είναι $10$ m και $13$ m.
\sol{12} Ελεύθερη κατασκευή. Το ύψος από το $\mathrm{Δ}$ είναι κάθετο στην $\mathrm{ΑΒ}$ και βγαίνει περίπου $2{,}6$ cm.
\sol{13} Οι $\hat{\mathrm{Α}}$ και $\hat{\mathrm{Δ}}$ είναι εντός και επί τα αυτά γωνίες των παραλλήλων $\mathrm{ΑΒ}$, $\mathrm{ΔΓ}$ με τέμνουσα την $\mathrm{ΑΔ}$. Οι εντός και επί τα αυτά γωνίες παραλλήλων είναι παραπληρωματικές, άρα $\hat{\mathrm{Α}}+\hat{\mathrm{Δ}}=180^\circ$.
\sol{14} Όχι. Σε παραλληλόγραμμο οι διαδοχικές γωνίες είναι παραπληρωματικές, αλλά $110^\circ+65^\circ=175^\circ\neq180^\circ$. Άρα οι $\mathrm{ΑΒ}$ και $\mathrm{ΔΓ}$ δεν είναι παράλληλες.
\sol{15} Με τέμνουσα την $\mathrm{ΔΗ}$ οι εντός και επί τα αυτά γωνίες $\hat{\mathrm{Δ}}$ και $\hat{\mathrm{Η}}$ έχουν άθροισμα $180^\circ$, άρα $\mathrm{ΔΕ}\parallel\mathrm{ΗΖ}$. Αν και οι $\mathrm{ΔΗ}$, $\mathrm{ΕΖ}$ ήταν παράλληλες, θα ήταν παραλληλόγραμμο και οι απέναντι πλευρές θα ήταν ίσες. Όμως $\mathrm{ΔΕ}=4\neq7=\mathrm{ΗΖ}$, άρα δεν είναι παραλληλόγραμμο. Είναι τραπέζιο με βάσεις $\mathrm{ΔΕ}$ και $\mathrm{ΗΖ}$.
\sol{16} Οι $\hat{\mathrm{Α}}$ και $\hat{\mathrm{Δ}}$ είναι παραπληρωματικές, άρα $\hat{\mathrm{Α}}=180^\circ-\hat{\mathrm{Δ}}$. Ομοίως $\hat{\mathrm{Β}}=180^\circ-\hat{\mathrm{Γ}}$. Επειδή $\hat{\mathrm{Δ}}=\hat{\mathrm{Γ}}$, είναι $\hat{\mathrm{Α}}=\hat{\mathrm{Β}}$.
\sol{17} Ελεύθερη απάντηση. Σε κάθε αντικείμενο πρέπει να βρεις αν υπάρχουν δύο ζεύγη παράλληλων πλευρών (παραλληλόγραμμο) ή ένα μόνο (τραπέζιο).
\sol{18} Ελεύθερη κατασκευή. $\hat{\mathrm{Γ}}=70^\circ$ και $\hat{\mathrm{Β}}=\hat{\mathrm{Δ}}=180^\circ-70^\circ=110^\circ$.
